\section{Chapter~\ref{sec:muscles}. Output to muscles}

The design of the output (motor units, the synaptic safety factor, recruitment by size, the muscle pair, the stretch loop, and rhythm generators) is measured directly; the stability condition for a loop with delay is a theorem; the internal model of the body is a well-founded hypothesis; the numbers in the figures are computed with the book's model.

{\footnotesize
\setlength{\tabcolsep}{3pt}\renewcommand{\arraystretch}{1.15}
\begin{longtable}{>{\raggedright}p{0.5cm}>{\raggedright}p{4.0cm}>{\raggedright}p{5.6cm}>{\raggedright}p{2.3cm}>{\raggedright\arraybackslash}p{1.7cm}}
\toprule
No. & Claim & Experiment and what was measured & Source & Status \\
\midrule\endfirsthead
\toprule
No. & Claim & Experiment and what was measured & Source & Status \\
\midrule\endhead
1 & Motor unit: one \nt{ACET} neuron controls a group of fibers, from several hundred to a couple of thousand; in muscles for fine adjustment the units are smaller & Human muscles, counts of motor axons and muscle fibers: on average about $340$ fibers per neuron in an interosseous muscle of the hand, about $410$ in the brachioradialis, about $1930$ in the medial gastrocnemius. The chapter gives no exact number for the smallest units. & \cite{Feinstein1955mu} & measured \\
2 & The synapse on a muscle releases several times more transmitter than is needed to make the fiber fire & Review of measurements at neuromuscular synapses: the safety factor (ratio of released transmitter to threshold) in adult mammals is usually $3$--$5$; in humans the margin rests more on the folds of the receiving membrane than on the size of the release. & \cite{WoodSlater2001} & measured \\
3 & The twitch~\eqref{eq:twitch} rises over a time $T_c$ of order $50\ms$ & Single motor units of the human hand during voluntary effort, force averaged over the unit's spikes: rise time $30$--$100\ms$, under $70\ms$ for $80\,\%$ of units. The function $(t/T_c)e^{1-t/T_c}$ is an approximation from a motor-unit pool model. & \cite{MilnerBrown1973rec, Fuglevand1993} & measured \\
4 & Sum of twitches~\eqref{eq:forcesum}: at $5$, $15$, and $40$ spikes/s the force is $0.2$--$1.1F_1$, $1.8$--$2.2F_1$, and about $5.4F_1$ (Fig.~\ref{fig:tetanus}) & Computed with \texttt{models/\allowbreak muscle\_\allowbreak models.py}, $T_c=50\ms$: $0.21$--$1.10$, $1.76$--$2.19$, and $5.28$--$5.49\,F_1$ over the last $0.3\,\text{s}$. Linear summation is a simplification: the model has no saturation of a real muscle. & \texttt{models/\allowbreak muscle\_\allowbreak models.py} & model \\
5 & Units are recruited by size; the weakest are a hundred times weaker than the strongest & Cat ventral roots: with increasing reflex input, the neurons with small spikes in the root, that is, the small ones, fire first. Human hand interosseous muscle: unit twitch forces range from $0.1$ to $10$~g and grow almost linearly with the force at which the unit is recruited. & \cite{Henneman1957, MilnerBrown1973rec} & measured \\
6 & The force step is proportional to the force, $\Delta F/F\approx q-1$~\eqref{eq:sizeprinciple}: floating point & Derived in the chapter from the geometric series of forces. Consistent with experiment: at large efforts, far fewer new units are recruited for the same force increment. & derived in the chapter; \cite{MilnerBrown1973rec} & theory \\
7 & The scatter of force grows in proportion to the force itself & Voluntary isometric effort of a thumb muscle: the standard deviation of force grows linearly with mean force; for the same effort evoked by electrical stimulation, the scatter is constant. Pool model: the linear growth comes from recruiting units in order of force. & \cite{Jones2002sdn} & measured \\
8 & The smoothness of movements is a consequence of signal-dependent noise & Model: the trajectory that minimizes end-point scatter under such noise reproduces measured eye and arm trajectories. There is no direct experiment that separates this cause from others. & \cite{HarrisWolpert1998} & hypothesis \\
9 & The difference of the forces of a muscle pair sets the torque, the sum sets the stiffness~\eqref{eq:antagonists} & Theory of impedance control by co-contraction. Human arm: the stiffness of each joint, measured with small perturbations, is roughly proportional to the torque at it; different co-contraction ratios change the shape and orientation of the hand stiffness ellipse. & \cite{Hogan1984, GomiOsu1998stiff} & measured \\
10 & The stretch sensor excites its own \nt{ACET} neuron and, through an inhibitory intermediary, switches off the opponent (Fig.~\ref{fig:antagonists}) & Cat spinal cord: a single intermediary excited by stretch-sensor fibers produces monosynaptic inhibitory potentials in the \nt{ACET} neurons of the opposing muscle, synaptic delay $0.28$--$0.42\ms$. The intermediary's transmitter (glycine) was not identified in this experiment. & \cite{Jankowska1972Ia} & measured \\
11 & Loop delay: spinal $25$--$30\ms$, through the cortex one and a half to three times longer, through the eye $100$--$200\ms$ & Human arm, a perturbation of the joint: short-latency muscle response at $20$--$45\ms$, long-latency at $45$--$105\ms$, voluntary at $120$--$180\ms$. A shift in the seen position of the finger during a movement: correction after $160\ms$ on average. & \cite{Pruszynski2008rapid, SaundersKnill2003vis} & measured \\
12 & $\dd x/\dd t=-Gx(t-d)$ is stable for $Gd<\pi/2$~\eqref{eq:delaystable}; at the boundary the frequency is $1/(4d)$ & Theorem on the roots of a delay equation. Check with \texttt{models/\allowbreak muscle\_\allowbreak models.py}, $d=30\ms$: by $3\,\text{s}$ the amplitude is $3\cdot10^{-6}$ at $G=40$, $0.13$ at $G=50$, $38$ at $G=55$ per second (boundary $52$). & \cite{Hayes1950} & theory \\
13 & Physiological tremor at $8$--$12$~Hz; the stretch loop is one of its sources & Review of measurements: fine tremor around $10$~Hz is present in healthy people; its origin is mixed (central oscillations, properties of unit discharges, mechanical resonance, and resonance of the reflex loop), and different sources dominate under different conditions. & \cite{McAuleyMarsden2000} & hypothesis \\
14 & The brain predicts the result of a command from an internal model of the body & A person holds a load in a pinch grip and moves the arm: under inertial, viscous, and elastic loads, grip force changes simultaneously with the load, not with the loop delay, so the load is predicted. A review collects such experiments; the model itself has not been observed directly in neurons. & \cite{FlanaganWing1997grip, WolpertGhahramani2000} & hypothesis \\
15 & The rhythm of walking, breathing, and chewing is generated by local networks under constant input & Review of experiments: an isolated nervous system (spinal cord, ganglia) produces a rhythmic motor pattern without rhythmic input and without feedback from muscles. & \cite{MarderBucher2001} & measured \\
16 & Mutual inhibition with fatigue~\eqref{eq:halfcentre} gives a rhythm with period $0.84\,\text{s}\approx2\tau_a$ (Fig.~\ref{fig:halfcentre}) & Theorem: a network with mutual inhibition and adaptation that has no stable state under constant input must oscillate. Computed with \texttt{models/\allowbreak muscle\_\allowbreak models.py} ($I=1$, $\beta=2$, $b=2$, $\tau=20\ms$, $\tau_a=0.4\,\text{s}$): period $0.84\,\text{s}$. That real generators are built exactly this way has not been shown. & \cite{Matsuoka1985osc} & model \\
\bottomrule
\end{longtable}}
